\begin{itemize}
    \item \textbf{背景と課題}:
    強相関電子系におけるモット転移を正確に記述する手法として,動的平均場理論が広く用いられている.

    しかしこの理論は,複雑な不純物問題を解くために極めて計算負荷の大きい厳密な数値計算を必要とする課題があった.

    \item \textbf{提案/目的}:
    本研究の目的は,動的平均場理論における無限個の電子の浴を少数のゴースト軌道で置き換える近似手法を用いて,モット転移を定量的かつ極めて低い計算コストで解析することである.

    \item \textbf{手法}:
    最も基礎的な強相関モデルである単一バンドハバードモデルを対象とした.

    物理軌道に加えて2つの空間ゴースト軌道を導入して変分空間を拡張し,自己無撞着な最適化計算を実装した.

    そして,無限次元Bethe格子および2次元正方格子における基底状態の性質を解析した.

    \item \textbf{主結果}:
    相互作用パラメータの増大に伴い,金属相の指標である準粒子重みがゼロとなる金属絶縁体転移を明確に観測した.

    さらに,昇順および降順のスイープ計算を行うことで,一次相転移に特有のヒステリシスループを捉えることに成功した.

    特に無限次元Bethe格子における転移点の臨界値は,動的平均場理論による厳密な計算結果と極めて良い一致を示した.

    \item \textbf{結論/意義}:
    浴の自由度を劇的に削減する大胆な近似を行ったにもかかわらず,本手法はモット転移の物理を高精度に記述できることが証明された.

    厳密対角化を用いて一瞬で解を導出できる圧倒的な計算コストの低さは,本手法が今後より複雑な現実の物質系へ適用可能な強力なツールとなることを示している.
\end{itemize}